\documentclass[10pt]{extarticle}
\usepackage[a4paper,margin=1.5cm]{geometry}
\usepackage[T1]{fontenc}
\usepackage{amsmath,amssymb}
\usepackage{graphicx}
\usepackage[export]{adjustbox}
\usepackage{multicol}
\usepackage{enumitem}
\usepackage{booktabs,array,tabularx}
\setlength{\columnsep}{0.7cm}
\setlength{\parindent}{0pt}
\setlength{\parskip}{3pt}
\pagestyle{empty}
\makeatletter
% figure registry (figdefs.tex is generated by the extract step)
\newcommand{\qfigdef}[4]{\@namedef{qf@file@#1}{#2}\@namedef{qf@w@#1}{#3}\@namedef{qf@wide@#1}{#4}}
\newcommand{\qf@check}[1]{\@ifundefined{qf@file@#1}{\PackageError{score2}{Unknown figure id #1}{}}{}}
\newcommand{\qf@img}[4][]{\includegraphics[width=#2\dimexpr\@nameuse{qf@w@#4}pt\relax,max width=#3,max height=0.45\textheight,#1]{\@nameuse{qf@file@#4}}}
% Full-width questions: finalize wraps every question that contains a wide figure in \qfullbegin...\qfullend,
% so the WHOLE question (stem, figures, options) is set across the page in reading order, then the two columns
% resume. The question's head - its number, stem and first figure (\qfullsplit marks the end of the head) -
% is kept on one page, so a number is never left at a page foot with its figure overleaf; the rest may break
% across pages normally, so no big gaps are left. Without a \qfullsplit, the whole question is kept together
% if it is under 45% of a page. Never write these yourself.
\newif\ifqf@full
\newif\ifqf@boxed
\newsavebox{\qf@box}
\newcommand{\qf@open}{\global\qf@fulltrue\par\global\qf@boxedtrue
  \setbox\qf@box\vbox\bgroup\hsize\textwidth\linewidth\textwidth}
\newcommand{\qf@emit}[1]{\par\egroup\global\qf@boxedfalse
  \ifdim\dimexpr\ht\qf@box+\dp\qf@box\relax<#1\textheight\noindent\usebox{\qf@box}\par\else\unvbox\qf@box\fi}
\newcommand{\qfullsplit}{\ifqf@boxed\qf@emit{0.85}\fi}
\newcommand{\qf@close}{\ifqf@boxed\qf@emit{0.45}\fi\global\qf@fullfalse}
\newcommand{\qfullbegin}{\end{multicols}\qf@open}
\newcommand{\qfullend}{\qf@close\begin{multicols}{2}}
\newcommand{\qfullstart}{\qf@open}
\newcommand{\qfulldone}{\qf@close}
% \qfig[scale]{id}: block figure, centred, directly where it is written (size set by the tool).
% Fallback if a wide figure is not inside a full-width question: break out of the columns in place - no floats.
\newcommand{\qfig}[2][1]{\qf@check{#2}%
  \ifqf@full\qf@blk{#1}{#2}\else
  \ifnum\@nameuse{qf@wide@#2}=1 \expandafter\@firstoftwo\else\expandafter\@secondoftwo\fi
  {\end{multicols}\begin{center}\qf@img{#1}{\textwidth}{#2}\end{center}\begin{multicols}{2}}%
  {\qf@blk{#1}{#2}}\fi}
\newcommand{\qf@blk}[2]{\par\smallskip{\centering\qf@img{#1}{\linewidth}{#2}\par}\smallskip}
% \qfigi[scale]{id}: inline figure (inside option lists / tables / minipages); never breaks columns.
\newcommand{\qfigi}[2][1]{\qf@check{#2}\qf@img[valign=c]{#1}{\linewidth}{#2}}
\makeatother
% structure macros - lint reads these to check completeness and figure placement
\newcommand{\papersec}[1]{\par\bigskip{\centering\Large\bfseries PAPER-#1\par}\medskip}
\newcommand{\partsec}[2]{\par\medskip\noindent{\large\bfseries PART-#1: #2}\par\smallskip}
\newcommand{\sect}[1]{\par\smallskip\noindent\textbf{#1}\par}
\newcommand{\qn}[1]{\par\smallskip\noindent\textbf{#1.}~}
\newcommand{\paraq}[2]{\par\smallskip\noindent\textbf{Paragraph for Questions #1 and #2}\par}
\newenvironment{qoptions}{\begin{itemize}[leftmargin=1.8em,itemsep=1pt,topsep=2pt,label={}]}{\end{itemize}}
\newcommand{\opt}[2]{\item[(#1)] #2}
% figure registry (generated by score2_tools.py finalize)
\qfigdef{p1_phy_s0_q01_1}{SRGBT - Test 07 (05 Sep 2021)-img/p1_phy_s0_q01_1.png}{105.4}{0}
\qfigdef{p1_phy_s0_q02_2}{SRGBT - Test 07 (05 Sep 2021)-img/p1_phy_s0_q02_2.png}{110.6}{0}
\qfigdef{p1_phy_s0_q02_1}{SRGBT - Test 07 (05 Sep 2021)-img/p1_phy_s0_q02_1.png}{112.8}{0}
\qfigdef{p1_phy_s0_q02_3}{SRGBT - Test 07 (05 Sep 2021)-img/p1_phy_s0_q02_3.png}{110.2}{0}
\qfigdef{p1_phy_s0_q02_4}{SRGBT - Test 07 (05 Sep 2021)-img/p1_phy_s0_q02_4.png}{113.6}{0}
\qfigdef{p1_phy_s0_q03_1}{SRGBT - Test 07 (05 Sep 2021)-img/p1_phy_s0_q03_1.png}{79.1}{0}
\qfigdef{p1_phy_s0_q03_2}{SRGBT - Test 07 (05 Sep 2021)-img/p1_phy_s0_q03_2.png}{121.9}{0}
\qfigdef{p1_phy_s0_q03_4}{SRGBT - Test 07 (05 Sep 2021)-img/p1_phy_s0_q03_4.png}{100.1}{0}
\qfigdef{p1_phy_s0_q03_5}{SRGBT - Test 07 (05 Sep 2021)-img/p1_phy_s0_q03_5.png}{90.0}{0}
\qfigdef{p1_phy_s0_q04_1}{SRGBT - Test 07 (05 Sep 2021)-img/p1_phy_s0_q04_1.png}{121.2}{0}
\qfigdef{p1_phy_s0_q05_1}{SRGBT - Test 07 (05 Sep 2021)-img/p1_phy_s0_q05_1.png}{130.1}{0}
\qfigdef{p1_phy_s0_q06_1}{SRGBT - Test 07 (05 Sep 2021)-img/p1_phy_s0_q06_1.png}{82.5}{0}
\qfigdef{p1_phy_s0_q07_1}{SRGBT - Test 07 (05 Sep 2021)-img/p1_phy_s0_q07_1.png}{97.4}{0}
\qfigdef{p1_phy_s0_q08_1}{SRGBT - Test 07 (05 Sep 2021)-img/p1_phy_s0_q08_1.png}{212.1}{0}
\qfigdef{p1_phy_s0_q09_1}{SRGBT - Test 07 (05 Sep 2021)-img/p1_phy_s0_q09_1.png}{88.1}{0}
\qfigdef{p1_phy_s0_q10_1}{SRGBT - Test 07 (05 Sep 2021)-img/p1_phy_s0_q10_1.png}{259.0}{1}
\qfigdef{p1_phy_s0_q11_1}{SRGBT - Test 07 (05 Sep 2021)-img/p1_phy_s0_q11_1.png}{106.5}{0}
\qfigdef{p1_phy_s0_q12_1}{SRGBT - Test 07 (05 Sep 2021)-img/p1_phy_s0_q12_1.png}{192.4}{0}
\qfigdef{p1_phy_s0_q12_2}{SRGBT - Test 07 (05 Sep 2021)-img/p1_phy_s0_q12_2.png}{82.6}{0}
\qfigdef{p1_phy_s0_q14_2}{SRGBT - Test 07 (05 Sep 2021)-img/p1_phy_s0_q14_2.png}{137.4}{0}
\qfigdef{p1_phy_s0_q17_1}{SRGBT - Test 07 (05 Sep 2021)-img/p1_phy_s0_q17_1.png}{121.8}{0}
\qfigdef{p1_phy_s0_q19_1}{SRGBT - Test 07 (05 Sep 2021)-img/p1_phy_s0_q19_1.png}{153.8}{0}
\qfigdef{p1_phy_s0_q20_1}{SRGBT - Test 07 (05 Sep 2021)-img/p1_phy_s0_q20_1.png}{92.3}{0}
\qfigdef{p1_chem_s0_q21_1}{SRGBT - Test 07 (05 Sep 2021)-img/p1_chem_s0_q21_1.png}{75.2}{0}
\qfigdef{p1_chem_s0_q22_1}{SRGBT - Test 07 (05 Sep 2021)-img/p1_chem_s0_q22_1.png}{245.2}{0}
\qfigdef{p1_chem_s0_q23_1}{SRGBT - Test 07 (05 Sep 2021)-img/p1_chem_s0_q23_1.png}{83.6}{0}
\qfigdef{p1_chem_s0_q23_3}{SRGBT - Test 07 (05 Sep 2021)-img/p1_chem_s0_q23_3.png}{111.7}{0}
\qfigdef{p1_chem_s0_q23_2}{SRGBT - Test 07 (05 Sep 2021)-img/p1_chem_s0_q23_2.png}{93.2}{0}
\qfigdef{p1_chem_s0_q23_5}{SRGBT - Test 07 (05 Sep 2021)-img/p1_chem_s0_q23_5.png}{111.5}{0}
\qfigdef{p1_chem_s0_q23_4}{SRGBT - Test 07 (05 Sep 2021)-img/p1_chem_s0_q23_4.png}{88.2}{0}
\qfigdef{p1_chem_s0_q25_1}{SRGBT - Test 07 (05 Sep 2021)-img/p1_chem_s0_q25_1.png}{51.8}{0}
\qfigdef{p1_chem_s0_q26_A}{SRGBT - Test 07 (05 Sep 2021)-img/p1_chem_s0_q26_A.png}{200.3}{0}
\qfigdef{p1_chem_s0_q26_2}{SRGBT - Test 07 (05 Sep 2021)-img/p1_chem_s0_q26_2.png}{182.7}{0}
\qfigdef{p1_chem_s0_q27_1}{SRGBT - Test 07 (05 Sep 2021)-img/p1_chem_s0_q27_1.png}{293.5}{1}
\qfigdef{p1_chem_s0_q28_1}{SRGBT - Test 07 (05 Sep 2021)-img/p1_chem_s0_q28_1.png}{328.0}{1}
\qfigdef{p1_chem_s0_q29_1}{SRGBT - Test 07 (05 Sep 2021)-img/p1_chem_s0_q29_1.png}{316.8}{1}
\qfigdef{p1_chem_s0_q30_1}{SRGBT - Test 07 (05 Sep 2021)-img/p1_chem_s0_q30_1.png}{245.2}{0}
\qfigdef{p1_chem_s0_q31_1}{SRGBT - Test 07 (05 Sep 2021)-img/p1_chem_s0_q31_1.png}{45.3}{0}
\qfigdef{p1_chem_s0_q31_2}{SRGBT - Test 07 (05 Sep 2021)-img/p1_chem_s0_q31_2.png}{38.0}{0}
\qfigdef{p1_chem_s0_q31_3}{SRGBT - Test 07 (05 Sep 2021)-img/p1_chem_s0_q31_3.png}{33.6}{0}
\qfigdef{p1_chem_s0_q31_4}{SRGBT - Test 07 (05 Sep 2021)-img/p1_chem_s0_q31_4.png}{42.6}{0}
\qfigdef{p1_chem_s0_q32_1}{SRGBT - Test 07 (05 Sep 2021)-img/p1_chem_s0_q32_1.png}{237.0}{0}
\qfigdef{p1_chem_s0_q32_3}{SRGBT - Test 07 (05 Sep 2021)-img/p1_chem_s0_q32_3.png}{276.5}{1}
\qfigdef{p1_chem_s0_q33_1}{SRGBT - Test 07 (05 Sep 2021)-img/p1_chem_s0_q33_1.png}{259.0}{1}
\qfigdef{p1_chem_s0_q34_1}{SRGBT - Test 07 (05 Sep 2021)-img/p1_chem_s0_q34_1.png}{245.2}{0}
\qfigdef{p1_chem_s0_q34_3}{SRGBT - Test 07 (05 Sep 2021)-img/p1_chem_s0_q34_3.png}{197.9}{0}
\qfigdef{p1_chem_s0_q38_1}{SRGBT - Test 07 (05 Sep 2021)-img/p1_chem_s0_q38_1.png}{117.2}{0}
\qfigdef{p1_chem_s0_q38_2}{SRGBT - Test 07 (05 Sep 2021)-img/p1_chem_s0_q38_2.png}{44.2}{0}
\raggedcolumns
\makeatletter\@beginparpenalty=10000\makeatother
\newcolumntype{L}[1]{>{\raggedright\arraybackslash}p{#1}}
\begin{document}
\begin{center}{\Large\bfseries SRGBT --- Test 07 (05 Sep 2021)}\end{center}
\begin{multicols}{2}
\raggedcolumns
\papersec{1}
\partsec{1}{PHYSICS}
\qn{1} A vessel having two immiscible fluids having density $\rho_1$ and $\rho_2$ filled upto height h$_1$ and h$_2$ placed on a block of height d as shown in figure. There is a hole at a distance x from the bottom of vessel. For all x $\in$ (0, h$_1$) if x increases, range (R) increases.\par
\qfig{p1_phy_s0_q01_1}
Choose the correct relation(s).
\begin{qoptions}
\opt{A}{$\dfrac{\rho_2\mathrm h_2}{\rho_1}>\mathrm d+\mathrm h_1$}
\opt{B}{$\dfrac{\rho_1\mathrm h_2}{\rho_2}>\mathrm d+\mathrm h_1$}
\opt{C}{$\dfrac{\rho_1\mathrm h_2}{\rho_2}>\mathrm d-\mathrm h_1$}
\opt{D}{$\dfrac{\rho_2\mathrm h_2}{\rho_1}>\mathrm d-\mathrm h_1$}
\end{qoptions}
\qn{2} With a wooden ball and a high water tank we hold four experiments. In the first experiment we have a floating ball weighed in a vessel. In the second experiment, the ball is weighted while tied to the bottom of the vessel, the third time the ball held under water by using the thin rod, finally, during the fourth experiment the ball is taken deep inside and released. In all the four cases weight is recorded. In which case the weight recorded is more than the weight of ball + water + tank?
\begin{qoptions}
\opt{A}{\qfigi{p1_phy_s0_q02_2}}
\opt{B}{\qfigi{p1_phy_s0_q02_1}}
\opt{C}{\qfigi{p1_phy_s0_q02_3}}
\opt{D}{\qfigi{p1_phy_s0_q02_4}}
\end{qoptions}
\qn{3} In table shown below \textbf{List- I} shows the possible conclusions regarding ``h'', for situation described in \textbf{List- II}. Match the possible outcomes.\par
\textbf{List- I}\par
(P)\par
\qfig{p1_phy_s0_q03_1}
Both containers are half filled with water \& arrangement is placed on weighing machine. Reading of weighing machine =``h''. Initially pump was off. The pump starts transmitting water from lower container to upper container. Then just after start of the pump :\par
(Q)\par
\qfig{p1_phy_s0_q03_2}
Stones are taken from container and gently released into water tank. Then :\par
(R)\par
\qfig{p1_phy_s0_q03_4}
Ice with wax embedded\par
When complete ice melts :\par
(S)\par
\qfig{p1_phy_s0_q03_5}
Closed at top by piston $\rho_\mathrm s$ = 0.65 g / cm$^3$\par
$\rho_\mathrm{air}$ = 1.30 kg / m$^3$\par
On increasing the pressure of air by moving piston downwards :\par
\textbf{List- II}\par
(1) h increases\par
(2) h decreases\par
(3) h remains same\par
(4) Cannot be perdicted from given information\par
\textbf{Codes :}\par
{\setlength{\tabcolsep}{8pt}\noindent\begin{tabular}{llllll} & P & Q & R & S\\ (A) & 1 & 3 & 2 & 4\\ (B) & 1 & 2 & 3 & 2\\ (C) & 3 & 2 & 4 & 1\\ (D) & 1 & 2 & 3 & 3\end{tabular}}\par
\qn{4} A glass is filled to a height h$_0$. Volume of liquid is V$_0$ and density is $\rho$ . A straw with uniform cross-sectional area `A' is used to drink liquid by creating a pressure at top P$_\mathrm{top}$ that is less than atmospheric pressure. Work done by the person in order to drink all the liquid in glass with constant velocity u (of liquid) is\par
\qfig{p1_phy_s0_q04_1}
\begin{qoptions}
\opt{A}{W = $\rho\mathrm V_0\left[\mathrm g\left(\mathrm h-\dfrac{\mathrm h_0}2\right)+\dfrac12\mathrm u^2\right]$}
\opt{B}{W = $\rho\mathrm V_0\left[\mathrm g(\mathrm h-\mathrm h_0)+\dfrac12\mathrm u^2\right]$}
\opt{C}{W = $\dfrac12\rho\mathrm V_0\mathrm u^2+\rho\mathrm V_0\mathrm{gh}$}
\opt{D}{W = $\dfrac12\rho\mathrm V_0\mathrm u^2+\rho\mathrm{gh}_0$}
\end{qoptions}
\qn{5} A liquid is placed in a cylindrical glass container of radius R$_0$. In gravity free space, free surface of liquid has spherical shape of radius of curvature R$_1$. Now same liquid is placed in a large spherical glass container of radius R$_2$ in uniform gravity. If surface of liquid in spherical container is found to be flat, then find the height of liquid surface in spherical container from centre of container :\par
\qfig{p1_phy_s0_q05_1}
\begin{qoptions}
\opt{A}{$\sqrt{\dfrac{\mathrm R_1\mathrm R_2}{\mathrm R_0}}$}
\opt{B}{$\dfrac{\mathrm R_1\mathrm R_2}{\mathrm R_0}$}
\opt{C}{$\mathrm R_0\dfrac{\mathrm R_1}{\mathrm R_2}$}
\opt{D}{$\mathrm R_0\dfrac{\mathrm R_2}{\mathrm R_1}$}
\end{qoptions}
\qn{6} A thin uniform elastic rod of natural length L, density $\rho$ and Young's modulus Y is just completely immersed in vertical position in fluid of density 2$\rho$ by applying a vertically downward force at its top end. The rod is in equilibrium, then pick up the correct option. (g is acceleration due to gravity; Neglect atmospheric pressure)\par
\qfig{p1_phy_s0_q06_1}
\begin{qoptions}
\opt{A}{There is net compression in rod of magnitude $\dfrac{\rho\mathrm L^2\mathrm g}{4\mathrm Y}$}
\opt{B}{There is net compression in rod of magnitude $\dfrac{3\rho\mathrm L^2\mathrm g}{2\mathrm Y}$}
\opt{C}{There is net compression in rod of magnitude $\dfrac{\rho\mathrm L^2\mathrm g}{2\mathrm Y}$}
\opt{D}{There is net compression in rod of magnitude $\dfrac{\rho\mathrm L^2\mathrm g}{\mathrm Y}$}
\end{qoptions}
\qn{7} A cylinder of radius R is kept embedded along the wall of A dam as shown. Take density of water as $\rho$. Take length as L.\par
\qfig{p1_phy_s0_q07_1}
Choose correct statement(s) :
\begin{qoptions}
\opt{A}{The vertical force exerted by water on the cylinder is $\dfrac{\rho\pi\mathrm R^2\mathrm{Lg}}2$}
\opt{B}{The net torque exerted by liquid on the cylinder is $\dfrac{\rho\mathrm R^3\mathrm{Lg}}2$}
\opt{C}{The force exerted by liquid on the cylinder in horizontal direction is [Neglecting atmospheric pressure] 3R$^2\rho$L}
\opt{D}{The force exerted by liquid on the cylinder in horizontal direction is [Neglecting atmospheric pressure] 2R$^2\rho$gL}
\end{qoptions}
\qn{8} A cubical vessel closed at top filled with a liquid of density $\rho$ is moving with acceleration g/2 along negative x-axis as shown in figure-1. The level of liquid in the vessel is as shown in figure. The side view of the level of liquid in vessel is as shown in the figure-2. The free part of vessel is evacuated :\par
\qfig{p1_phy_s0_q08_1}
\begin{qoptions}
\opt{A}{Pressure at point (x, z) inside vessel will be non-zero for z $>\dfrac{\mathrm x}2+(\mathrm L-\ell)$}
\opt{B}{Pressure at point (x, z) inside vessel will be $(\mathrm L-\ell+\mathrm x-\mathrm z)\rho\mathrm g$ for z $>\dfrac{\mathrm x}2+(\mathrm L-\ell)$}
\opt{C}{Pressure at point (x, z) inside vessel will be 0 for z $<\dfrac{\mathrm x}2+(\mathrm L-\ell)$}
\opt{D}{Pressure at point (x, z) inside vessel will be $\left(\mathrm L-\ell+\dfrac{\mathrm x}2-\mathrm z\right)\rho\mathrm g$ for z $<\dfrac{\mathrm x}2+(\mathrm L-\ell)$}
\end{qoptions}
\qn{9} A spherical ball of radius R is made by joining two hemispherical parts. The two parts have density $\rho$ and 2$\rho$. When placed in a water tank, the ball floats while remaining completely submerged.\par
\qfig{p1_phy_s0_q09_1}
Neglect viscous forces. Mark the correct statement(s) :
\begin{qoptions}
\opt{A}{Density of water is $\rho_0$, then $\rho=\dfrac{2\rho_0}3$.}
\opt{B}{Position of centre of mass of ball is at distance of $\dfrac{\mathrm R}8$ from O.}
\opt{C}{Moment of inertia of ball about COM of ball perpendicular to the plane of paper is $\dfrac{123}{160}\rho\pi\mathrm R^5$}
\opt{D}{Time period of small angular oscillations of the ball about its equilibrium position is given by $2\pi\sqrt{\dfrac{123\mathrm R}{40\mathrm g}}$}
\end{qoptions}
\qfullbegin
\qn{10} Two identical cylinders with pistons of area S$_2$ are connected by a pipe. In the cylinder is water. On top of the pistons we put identical cylindrical cups of area S$_1$ with equal amounts of water. Then in one of the cups we lower a body of mass m, and another a body of mass M. Neither the first body nor second sinks to the bottom. The system is allowed to come to equilibrium as shown in figure 2 (M $>$ m)\par
\qfig{p1_phy_s0_q10_1}
\qfullsplit
\begin{qoptions}
\opt{A}{Value of x is $\dfrac{\mathrm M-\mathrm m}{(\mathrm S_1-\mathrm S_2)\rho}$}
\opt{B}{Value of x is $\dfrac{(\mathrm M-\mathrm m)(\mathrm S_2-\mathrm S_1)}{\mathrm S_1\mathrm S_2\rho}$}
\opt{C}{Value of y is $\dfrac{(\mathrm M-\mathrm m)}{\mathrm S_2\rho}$}
\opt{D}{Value of y is $\dfrac{(\mathrm M-\mathrm m)}{\mathrm S_1\rho}$}
\end{qoptions}
\qfullend
\qn{11} In the figure, AB and CD are two vertical walls of a horizontal rectangular groove. A rectangular block of mass m just fits in the groove. Base of the block is smooth and coefficient of friction between walls and block is $\mu$. Contact area of block with each wall is A. Linear coefficient of thermal expansion and young's modulus of material of block are $\alpha$ and Y respectively. Temperature $\theta$ of block increases uniformly as $\Delta\theta$ = t (t = time). Initial velocity of block is v$_0$ then\par
\qfig{p1_phy_s0_q11_1}
\begin{qoptions}
\opt{A}{Retardation of block at time t is $\dfrac{\mu\mathrm{YA}\alpha\mathrm t}{\mathrm m}$}
\opt{B}{Velocity at time t is $\left(\mathrm v_0-\dfrac{\mu\mathrm{YA}\alpha\mathrm t^2}{\mathrm m}\right)$}
\opt{C}{Time taken by block to stop is $\sqrt{\dfrac{\mathrm{mv}_0}{\mu\mathrm{YA}\alpha}}$}
\opt{D}{work done by friction is $\dfrac12\mathrm{mv}_0^2$}
\end{qoptions}
\qn{12} In the figure shown, is given a flat plate and a rotating cone with a very acute angle $\theta$ (less than 5$^\circ$) The apex of cone just touches the plate surface. A liquid with coefficient of viscosity $\eta$ fills the narrow gap formed by the cone and plate. Cone is rotated with angular velocity $\omega$ as shown. Neglect end effects and assume a linear velocity profile :\par
\qfig{p1_phy_s0_q12_1}
\begin{qoptions}
\opt{A}{Velocity gradient in the liquid is $\dfrac\omega\theta$.}
\opt{B}{Velocity gradient in the liquid is $\theta$.}
\opt{C}{The torque on the driven cone is equal to $\left(\dfrac23\pi\mathrm R^3\eta\right)\dfrac\omega\theta$}
\opt{D}{The torque on the driven cone is equal to $\left(\dfrac13\pi\mathrm R^3\eta\right)\dfrac\omega\theta$}
\end{qoptions}
\paraq{13}{14}Viscosity is measured by a device called the rotating-cylinder viscometer which is shown schematically in figure. The outer cylinder, which has a bottom so as to contain the liquid, is mounted on a variable-speed turntable. An inner cylinder hangs from a torsion balance calibrated so that its scale directly reads the torque exerted on the inner cylinders. The two cylinders are coaxial. The difference between their radii, r$_2$ $-$ r$_1$, is much smaller than either radius. Thus the annular (ring-shaped) space between the two cylinders is a reasonable approximation of the planar space between the two parallel plates of figure.\par
The space between the cylinders is filled to a depth h with the fluid whose viscosity is to be measured. The outer cylinder is made to rotate, and it pulls the fluid along with it. As a result of the viscous drag of the fluid, the inner cylinder experiences a torque $\tau$. Torsional stiffness of torsion fibre is C.\par
\qfig{p1_phy_s0_q12_2}
\qn{13} What is average velocity gradient in space between cylinders ?
\begin{qoptions}
\opt{A}{$\dfrac{\omega(\mathrm r_2-\mathrm r_1)}{\mathrm h}$}
\opt{B}{$\dfrac{\omega\mathrm r_1}{\mathrm h}$}
\opt{C}{$\dfrac{\omega\mathrm r_2}{\mathrm h}$}
\opt{D}{$\dfrac{\omega\mathrm r_2}{(\mathrm r_2-\mathrm r_1)}$}
\end{qoptions}
\qn{14} Find angle through which torsional fibre gets twisted if turn table is made to reach constant $\omega$ slowly.
\begin{qoptions}
\opt{A}{$\theta=\dfrac{2\pi\eta\omega\mathrm h(\mathrm r_2-\mathrm r_1)}{\mathrm{Cr}_1\mathrm r_2}$}
\opt{B}{$\theta=\dfrac{2\pi\eta\omega\mathrm r_1\mathrm r_2}{\mathrm{Ch}}$}
\opt{C}{$\theta=\dfrac{2\pi\eta\omega\mathrm{hr}_1^2\mathrm r_2}{\mathrm C(\mathrm r_2-\mathrm r_1)}$}
\opt{D}{$\theta=\dfrac{2\pi\eta\omega(\mathrm r_2-\mathrm r_1)}{\mathrm{Chr}_1\mathrm r_2}$}
\end{qoptions}
\paraq{15}{16}The fact that the density of ice is smaller than the density of water has an important consequence for homeowners, who have to contend with the possibility of bursting water pipes during severe winters. Water often freezes in a section of pipe exposed to unusually cold temperatures. The ice can form an immovable plug that prevents the subsequent flow of water, as Figure illustrates. When water (larger density) turns to ice (smaller density), its volume expands by 10\%. When ice forms on the right side of the plug, the expanding ice pushes liquid to the right. But it has nowhere to go if the faucet is closed. As ice continues to form and expand, the water pressure between the plug and faucet rises. The bulk modulus of ice is much more than that of the water. So water gets compressed. Even a small increase in the amount of ice produces a large increase in the pressure. Therefore, the pipe can burst at any point where it is structurally weak. There is a simple way to prevent pipes from bursting. Simply open the faucet, so it drips a little. The excessive pressure will be relieved.\par
\qfig{p1_phy_s0_q14_2}
\qn{15} Assume that the plug of ice formed expands only to the right and the mass of ice formed is equal to the mass of water to the right of it. If the bulk modulus of water is B = 2.2 $\times$ 10$^9$ Pa, What is the excess pressure of the water enclosed.(Atmospheric pressure is 1 $\times$ 10$^5$ Pa).
\begin{qoptions}
\opt{A}{2.2 $\times$ 10$^7$ Pa}
\opt{B}{2.2 $\times$ 10$^8$ Pa}
\opt{C}{2.2 $\times$ 10$^9$ Pa}
\opt{D}{2.2 $\times$ 10$^{10}$ Pa}
\end{qoptions}
\qn{16} Which of the following steps would make the pipe more capable of withstanding the increased pressure?
\begin{qoptions}
\opt{A}{Increase the radius of pipe and decrease thickness of it's walls.}
\opt{B}{Increase the radius of pipe and thickness of it's walls.}
\opt{C}{decrease the radius of pipe and thickness of it's walls.}
\opt{D}{decrease the radius of the pipe and increase the thickness of its walls.}
\end{qoptions}
\qn{17} A hollow thin walled closed cylindrical container can move inside a long cylindrical tube filled with a Newtonian viscous liquid. There is an extremely small gap between outer boundary of cylinder and inner boundary of tube. If cylinder is completely empty and released in the tube it acquires a terminal speed 6 cm/s in upward direction. When cylindrical container is completely filled with the same Newtonian liquid and then released it acquires a terminal speed of 2 cm/s in the downward direction. If cylindrical container is half filled with same newtonial liquid, it acquires a terminal speed of v in upward direction. Find v (in cm/s).\par
\qfig{p1_phy_s0_q17_1}
\qn{18} A closed vertical cylinder having a radius R and height equal to radius R is filled with oil of density $\rho$. If the cylinder is rotated about its axis at an angular velocity $\sqrt{\dfrac{\mathrm g}{\mathrm R}}$, the thrust at the bottom of the cylinder is $\left(\dfrac{\mathrm n}{\mathrm m}\right)\pi\rho\mathrm{gH}^3$, then \textbf{n+m} equals
\qn{19} A rod AB of mass m, and length $\ell$ is moving with constant acceleration a. The graph represents the stress develop at any cross-section versus position of cross-section from end A of rod. The density of rod is not uniform. If the density of rod at point-2 is $25\times10^3\,$kg/m$^3$ and that of at point-1 is k$\times10^3\,$kg/m$^3$ . Find the value of k.\par
\qfig{p1_phy_s0_q19_1}
\qn{20} A jet of liquid is split into two liquid jets after colliding with an incline plane, find the ratio of rate of flow of liquid jets just after collision. Assume after collision flow remains streamline and no losses of energy due to collision.\par
\qfig{p1_phy_s0_q20_1}
\raggedcolumns
\papersec{1}
\partsec{2}{CHEMISTRY}
\qn{21} Which of the following C --C bond length is minimum in anthracene?\par
\qfig{p1_chem_s0_q21_1}
\begin{qoptions}
\opt{A}{4}
\opt{B}{3}
\opt{C}{2}
\opt{D}{1}
\end{qoptions}
\qn{22} Which of the following order is incorrect ?\par
\qfig{p1_chem_s0_q22_1}
\qn{23} The most stable conformation of given compound.\par
\qfig{p1_chem_s0_q23_1}
\begin{qoptions}
\opt{A}{\qfigi{p1_chem_s0_q23_3}}
\opt{B}{\qfigi{p1_chem_s0_q23_2}}
\opt{C}{\qfigi{p1_chem_s0_q23_5}}
\opt{D}{\qfigi{p1_chem_s0_q23_4}}
\end{qoptions}
\qn{24} Which of the following statement is correct ?
\begin{qoptions}
\opt{A}{Nitrobenzene will give meta--nitrotoluene on reaction with CH$_3$Cl/AlCl$_3$.}
\opt{B}{Chlorobenzene will give meta--substituted product on electrophilic substitution since it exerts --I $>$ +M effect.}
\opt{C}{n--Propyl benzene can be easily obtained on Friedal crafts alkylation of benzene with n--propyl chloride.}
\opt{D}{Toluene can be obtained in better yield when excess of benzene will react with CH$_3$Cl/AlCl$_3$.}
\end{qoptions}
\qn{25}
\qfig{p1_chem_s0_q25_1}
$\xrightarrow{\Delta}$ $\xrightarrow{\text{AlCl}_3/\text{PhH}}$ $\xrightarrow{\text{Zn--Hg/HCl}}$ $\xrightarrow{\text{SOCl}_2}$ $\xrightarrow[\text{anhydrous}]{\text{AlCl}_3}$ $\xrightarrow{\text{Zn--Hg/HCl}}$ $\xrightarrow{\text{Pd}/\Delta}$\par
Final product is
\begin{qoptions}
\opt{A}{Naphthalene}
\opt{B}{Anthracene}
\opt{C}{Phenanthrene}
\opt{D}{Biphenyl}
\end{qoptions}
\qn{26} Which of the following acid base reaction is not fesible.
\begin{qoptions}
\opt{A}{\qfigi{p1_chem_s0_q26_A}}
\opt{B}{Me $-$ C $\equiv$ CH + NaOH $\to$ MeC $\equiv\overset{-}{\mathrm C}\overset{\oplus}{\mathrm{Na}}$ + H$_2$O}
\opt{C}{\qfigi{p1_chem_s0_q26_2}}
\opt{D}{Me$_3$CONa + H$_2$O $\to$ Me$_3$COH + NaOH}
\end{qoptions}
\qfullbegin
\qn{27} Which one of the following is least basic in character ?\par
\qfig{p1_chem_s0_q27_1}
\qn{28}
\qfig{p1_chem_s0_q28_1}
\qn{29} Which of the given following stability order is/are correct.\par
\qfig{p1_chem_s0_q29_1}
\qfullend
\qn{30} Correct increasing order of acidic strength of following compounds is / are ?\par
\qfig{p1_chem_s0_q30_1}
\qn{31} Which of the following will give only those products on ozonolysis which are obtained by reductive ozonolysis of m-xylene
\begin{qoptions}
\opt{A}{\qfigi{p1_chem_s0_q31_1}}
\opt{B}{\qfigi{p1_chem_s0_q31_2}}
\opt{C}{\qfigi{p1_chem_s0_q31_3}}
\opt{D}{\qfigi{p1_chem_s0_q31_4}}
\end{qoptions}
\qn{32} Observe the following reaction and select true statement about it's mechanism.\par
\qfig{p1_chem_s0_q32_1}
\begin{qoptions}
\opt{A}{It is an electrophilic substitution reaction}
\opt{B}{The leaving group is (CH$_3$)$_3$C$^\oplus$ ion}
\opt{C}{The H$^\oplus$ ion of HBr has been eliminated from (CH$_3$)$_3$C$^\oplus$ ion}
\opt{D}{It is an example of IPSO substitution reaction.}
\end{qoptions}
\qfullbegin
\paraq{33}{34}Observe the following reaction and answer the following questions.\par
\qfig{p1_chem_s0_q32_3}
\qn{33} The product `R' is :\par
\qfig{p1_chem_s0_q33_1}
\qfullend
\qn{34} The structure of Q$_1$ is\par
\qfig{p1_chem_s0_q34_1}
\paraq{35}{36}Consider the following figure representing variation of heat of combustion of isomeric butenes.\par
\qfig{p1_chem_s0_q34_3}
Based on this answer the following questions.
\qn{35} Minimum dipole moment is of isomers.
\begin{qoptions}
\opt{A}{P}
\opt{B}{Q}
\opt{C}{R}
\opt{D}{S}
\end{qoptions}
\qn{36} Which gives fastest reaction with HBr?
\begin{qoptions}
\opt{A}{P}
\opt{B}{Q}
\opt{C}{R}
\opt{D}{S}
\end{qoptions}
\qn{37} How many saturated bicyclo hydrocarbon are possible having molecular formula C$_8$H$_{14}$ with no primary hydrogen and no quaternary carbon?
\qn{38} The number of resoance structures for Z is:\par
\qfig{p1_chem_s0_q38_1}
\qn{39} The number of deuterium atoms found in the tautomer of \qfigi{p1_chem_s0_q38_2} when it is kept in OD$^\ominus$ / D$_2$O for a long period of time?
\qn{40} Consider the following groups :\par
$-$NO, $-\overset{+}{\mathrm S}\mathrm R_2$, $-$NO$_2$, $-$CMe$_3$, $-$OCOCH$_3$, $-\mathrm C(=\mathrm O)-$NHCH$_3$, $-\overset{+}{\mathrm N}$H$_3$, $-\overset{\ominus}{\mathrm B}$F$_3$, $-$CH$_2$--CH$_3$, $-$OH, $-$SH\par
Among given compounds.\par
(a) Number of groups which can show $-$M effect = P\par
(b) Number of groups which can show +M effect = Q\par
(c) Number of group which can show +H effect = R\par
\textbf{Give your answer in 3 digits as PQR.}
\raggedcolumns
\papersec{1}
\partsec{3}{MATHEMATICS}
\qn{41} Let a variable chord PQ of parabola $y^2=8x$ always passes through point (6, 0). If OPRQ is a parallelogram, then locus of point R is (O being the origin) :
\begin{qoptions}
\opt{A}{$y^2=4(x-6)$}
\opt{B}{$y^2=8(x-6)$}
\opt{C}{$y^2=4(x-12)$}
\opt{D}{$y^2=8(x-12)$}
\end{qoptions}
\qn{42} The locus of feet of perpendicular from either foci of the ellipse $(x-y+1)^2+(2x+2y-6)^2=20$ on any tangent will be :
\begin{qoptions}
\opt{A}{$x^2+y^2+2x+4y-20=0$}
\opt{B}{$x^2+y^2+2x+4y-5=0$}
\opt{C}{$x^2+y^2-2x-4y-5=0$}
\opt{D}{$x^2+y^2-2x-4y-20=0$}
\end{qoptions}
\qn{43} The product of the lengths of perpendiculars drawn from any point on the hyperbola $x^2-2y^2-2=0$ to its asymptotes is :
\begin{qoptions}
\opt{A}{$\dfrac12$}
\opt{B}{$\dfrac23$}
\opt{C}{$\dfrac32$}
\opt{D}{1}
\end{qoptions}
\qn{44} The equation of the line touching both the parabolas $y^2=4x$ and $x^2=-32y$ is :
\begin{qoptions}
\opt{A}{$x+2y+4=0$}
\opt{B}{$2x+y+4=0$}
\opt{C}{$x-2y-4=0$}
\opt{D}{$x-2y+4=0$}
\end{qoptions}
\qn{45} Let the normal to a parabola $y^2=4x$ makes an angle $\tan^{-1}2$ with the axis of parabola. If it will cut the curve again at angle $\alpha$, then $\alpha$ is-
\begin{qoptions}
\opt{A}{$\dfrac\pi2$}
\opt{B}{$\dfrac\pi4$}
\opt{C}{$\tan^{-1}2$}
\opt{D}{$\tan^{-1}\left(\dfrac12\right)$}
\end{qoptions}
\qn{46} If a tangent to ellipse $x^2+4y^2=4$, meet the ellipse $x^2+2y^2=6$ at points P \& Q, then angle between tangents at P and Q is-
\begin{qoptions}
\opt{A}{$30^\circ$}
\opt{B}{$60^\circ$}
\opt{C}{$75^\circ$}
\opt{D}{$90^\circ$}
\end{qoptions}
\qn{47} An equilateral triangle is inscribed in ellipse whose equation is $x^2+4y^2=4$. One vertex of triangle is (0,1) and one altitude is contained in y-axis \& length of each side is $\sqrt{\dfrac mn}$ (where m \& n are relatively prime), then m + n is :
\begin{qoptions}
\opt{A}{937}
\opt{B}{973}
\opt{C}{793}
\opt{D}{739}
\end{qoptions}
\qn{48} Locus of centre of ellipse sliding along co-ordinate axis is :
\begin{qoptions}
\opt{A}{straight line}
\opt{B}{circle}
\opt{C}{parabola}
\opt{D}{another ellipse}
\end{qoptions}
\qn{49} The parabola $y=x^2+px+q$ cuts the straight line $y=2x-3$ at a point with abscissa 1. If the distance between the vertex of the parabola and the x $-$ axis is least then :
\begin{qoptions}
\opt{A}{p = 0 \& q = $-2$}
\opt{B}{p = $-2$ \& q = 0}
\opt{C}{least distance between the parabola and x $-$ axis is 2}
\opt{D}{least distance between the parabola and x $-$ axis is 1}
\end{qoptions}
\qn{50} Maximum distance of any point on the ellipse $x^2+2y^2+2xy-1=0$ from the centre is-
\begin{qoptions}
\opt{A}{$\dfrac2{3-\sqrt5}$}
\opt{B}{$\dfrac{\sqrt5-1}2$}
\opt{C}{$\dfrac{\sqrt2}{\sqrt{3-\sqrt5}}$}
\opt{D}{$\dfrac{\sqrt5+1}2$}
\end{qoptions}
\qn{51} A tangent is drawn at any point P(t) on the parabola $y^2=8x$ and a point Q($\alpha,\beta$) is taken on it from which pair of tangents QA and QB are drawn to the circle $x^2+y^2=4$, where A \& B are points of contact. Now if point P(t) varies on the parabola then the locus of all such point of concurrency of the chord of contact AB is denoted by S = 0, then-
\begin{qoptions}
\opt{A}{Equation of S = 0 is $2x+y^2=0$}
\opt{B}{Equation of S = 0 is $x+y^2=0$}
\opt{C}{Area bounded by S = 0 and the line $2x+1=0$ is $\dfrac23$ sq. units}
\opt{D}{Tangents drawn at the extremities of intersection of S = 0 and the line $2x-y+1=0$ to the curve S = 0 includes an angle $75^\circ$.}
\end{qoptions}
\qn{52} Consider the ellipse $\dfrac{x^2}4+\dfrac{y^2}{b^2}=1$ (b$^2<4$) and whose eccentricity is `e'. If the normal at the end of latus-rectum passes through one extremity of minor axis, then which of the following statement(s) is/are true?
\begin{qoptions}
\opt{A}{Equation of director circle of ellipse is $x^2+y^2=10-2\sqrt5$}
\opt{B}{Length of latus rectum of the ellipse is equal to b$^2$}
\opt{C}{The value of e$^2$ + e is 1}
\opt{D}{Area of region between director and auxillary circle of the ellipse is $2\pi\left(3-\sqrt5\right)$}
\end{qoptions}
\paraq{53}{54}Let S$_1$ : $12x^2-7xy-12y^2+25=0$ is a hyperbola, L : y = 2x is a line and S$_2$ is reflection of S$_1$ in L.
\qn{53} Equation of S$_2$ is xy = c, then c is -
\begin{qoptions}
\opt{A}{$-2$}
\opt{B}{$-1$}
\opt{C}{1}
\opt{D}{2}
\end{qoptions}
\qn{54} Equation of transverse axis of S$_1$ is -
\begin{qoptions}
\opt{A}{y = 7x}
\opt{B}{x + 7y = 0}
\opt{C}{3x = 4y}
\opt{D}{7x + y = 0}
\end{qoptions}
\paraq{55}{56}Through the vertex A of parabola $y^2=4ax$, chord AP is drawn and chord PP$'$ is drawn at right angle to AP to meet the axis in Q.
\qn{55} Projection of PQ on x-axis is equal to-
\begin{qoptions}
\opt{A}{4a}
\opt{B}{2a}
\opt{C}{16a$^2$}
\opt{D}{4a$^2$}
\end{qoptions}
\qn{56} Locus of mid point of PP$'$ is-
\begin{qoptions}
\opt{A}{$2axy=(y^2+4a^2)^2+16a^4$}
\opt{B}{$2axy^2=(y^2+4a^2)^2+16a^4$}
\opt{C}{$2ax^2y=(y^2+4a^2)^2+16a^4$}
\opt{D}{$2ax^2y^2=(y^2+4a^2)^2+16a^4$}
\end{qoptions}
\qn{57} For each real number m the parabola $y=(m^2+4)x^2+(m-2)^2x-4m+2$ passes through the same point (a, b), then $a^2+b^2$ is equal to
\qn{58} The locus of mid point of those chords of hyperbola xy = 6 which intersects co-ordinate axes at A \& B such that radius of circumcircle of $\Delta$OAB (O is origin) is 5, is $x^2+y^2=a^2$. The value of $|$a$|$ is
\qn{59} Let ($-$1,$-$1) be the focus of a parabola and 3x $-$ y = 8 be the tangent drawn to it at point (7 , 13) on the parabola. If ax + by + 19 = 0 is the equation of directrix of the parabola then the value of a+b is (where a, b $\in$ N)
\qn{60} If P is a variable point on the curve $(x-y)^2=8(x+y-2)$ such that the tangents from point P to the curve $\dfrac{x^2}2+\dfrac{y^2}1=1$ are at right angles, then number of such points is
\end{multicols}
\end{document}
